%% 302_Elementarzellen_Spektrale_Antwort_En_ch.tex
%% FFGFT -- Doc. 302: Elementary cells, one-bit capacity, and partition invariance;
%%                    a spectral answer (information unit = winding quantum Delta w = 1)
%% Author: Johann Pascher
%% Date: July 2026
%% Language: English

\section*{Doc.~302 \quad Elementary cells, one-bit capacity, and partition invariance ---
a spectral answer}

The following addresses a natural objection to any \enquote{quantised space} picture:
what observer-independent criterion distinguishes one elementary cell from its
neighbours; why a fundamental length should imply fixed cells each storing exactly
one bit; whether the resulting entropy is invariant under changes of origin,
orientation, coordinates, or foliation; and whether the indivisibility of a cell
follows from the Planck scale or is an added assumption. The short answer is that
every one of these points is correct \emph{against a spatial-cell ontology}, and that
FFGFT does not rest on one. It replaces cells with topological and spectral invariants (winding classes and modes), and the objections
dissolve in turn.

\section{The picture the objections target}

A \enquote{one bit per Planck cell} ontology tiles space (or a horizon area) into
fixed cells of size ${\sim}\,\ell_P$, assigns each a one-bit capacity, and reads
entropy off the count of cells. Against that picture the following points are
decisive:

\begin{itemize}
\item there is no observer-independent criterion individuating one spatial cell from
its neighbour: a tiling carries a chosen origin and orientation;
\item the Planck length is a \emph{scale}, and a scale does not by itself imply a
partition into discrete containers;
\item a partition-based entropy is not invariant under change of origin, orientation,
coordinates, or foliation --- it inherits the arbitrariness of the partition;
\item if space is a continuum, any cell subdivides, so \enquote{one cell $=$ one
indivisible bit} needs an independent reason why half-cell states are forbidden;
\item and absent that reason, the cell is simply an extra postulate riding on top of
the Planck scale.
\end{itemize}

All five are granted. This is precisely why the framework is not built on spatial
cells.

\section{Invariants, not cells}

The fundamental object is not a lattice of cells but a compact, boundaryless manifold
--- a four-torus $T^4$ (with a $\mathbb{Z}_3$ identification), carried into a Hilbert
space $\mathcal{H}=L^2(T^4)\otimes\mathbb{C}^3$. The elementary discrete objects of
the theory are not regions of space; they are \emph{winding classes}
($\pi_1(T^4)\cong\mathbb{Z}^4$, the carriers of the information unit) and
\emph{eigenmodes} of an operator on that manifold (the connection-Laplacian / mass
operator, the carriers of the mass spectrum). With that single change of ontology ---
invariants, not cells --- the points resolve in turn.

\begin{enumerate}
\item \textbf{Observer-independent individuation.} The elementary discrete objects are
distinguished by \emph{invariants}, not by a location in space: the information unit by
its \emph{winding number} (a topological invariant, $\pi_1(T^4)$), the mass spectrum by
\emph{eigenvalues} (a spectral invariant). Neither depends on coordinates, origin, or
orientation. There is no \enquote{neighbouring cell in space} to distinguish; there are
discrete classes labelled by invariants. The individuation is topological or spectral,
not spatial --- which is exactly what makes it observer-independent.

\item \textbf{What the information unit is --- and what stays a definition in the
\enquote{one bit}.} The discreteness does not come from the Planck length and is not a
tiling; it is \emph{topological}: $\pi_1(T^4) \cong \mathbb{Z}^4$ admits only integer
winding numbers. The minimal information unit is therefore \emph{by definition} the
winding quantum $\Delta w = 1$ (Doc.~257) --- scale-universal and energy-independent;
there is no half winding quantum. The identification \enquote{one bit
$\;\widehat{=}\;$ one minimal winding} is a \emph{definition}, though a topologically
motivated one (the smallest non-trivial winding), not an arbitrary two-state
assumption. The $\ln 2$ is \emph{not} the fundamental unit but its thermodynamic
shadow: the \emph{energy} of a bit is scale-specific, $E_\text{bit}(L) = \hbar c / L$,
and Landauer's $E_L = k_B T \ln 2$ is the special case at a single (the biological)
scale. This answers Takayuki's point~(2): \enquote{why exactly one bit} is a definition
on a topological basis ($\Delta w = 1$), not forced by a scale --- and the unit is
scale-universal, hence precisely not tied to a chosen partition.

\item \textbf{Invariance under origin, orientation, coordinates, foliation.} Yes.
Because the entropy is a function of the spectrum --- a count or trace over eigenvalues
--- and eigenvalues are invariants, the entropy is automatically independent of origin,
orientation, coordinate choice, and basis. This is not an added axiom; it is a
property of spectral quantities. A partition-based entropy would fail the invariance
test, which is one reason the framework does not define entropy through a partition at
all. A further, framework-specific result makes the origin/orientation independence
explicit: on the compact structure there is no preferred index assignment --- all
$D_6$-equivalent segments of the torus structure are equally valid descriptions of the
same object (\enquote{meshing, not a cut-out}). No labelling of a \enquote{first} cell
or a preferred axis is available, or needed.

\item \textbf{Continuous space, no half-cells.} Space is continuous: $T^4$ is a smooth
manifold, with no spatial atomism. One therefore cannot obtain indivisibility by
subdividing space, and the framework does not try to. The indivisibility is
\emph{topological and spectral}, not spatial: $\pi_1(T^4)\cong\mathbb{Z}^4$ admits no
half winding, and the spectrum on the compact manifold has a smallest gap fixed by the
topology (the circumference / boundarylessness), just as a finite string
cannot vibrate at \enquote{half its fundamental}. \enquote{Half-cell states} is, in
this language, a category error: there are no cells to halve. The continuum of the
manifold and the discreteness of the spectrum coexist without tension --- exactly as in
acoustics, where a continuous medium supports a discrete set of normal modes.

\item \textbf{Derived, not assumed.} The discreteness follows; it is not an extra
\enquote{cell postulate}. Given a compact manifold, the discreteness of the spectrum
is a theorem. What \emph{is} assumed is the compact geometry itself --- the torus
$T^4/\mathbb{Z}_3$ and the value of $\xi$. So the honest accounting: the
\enquote{one-bit-per-cell} postulate is traded for a \enquote{compact-geometry}
postulate, and from the latter the discreteness is derived rather than injected. That
is a real assumption, and it is where the framework should be attacked --- but
empirically, not by fiat: the same geometry that yields the discrete spectrum also
fixes particle masses, and those are falsifiable numbers (for instance a $\tau$ mass
of $1776.969~\mathrm{MeV}$, checkable at Belle~II). If the geometry is wrong, the
spectrum and the masses fail together.
\end{enumerate}

\section{What is assumed, derived, and posited}

The geometry itself is also a \emph{description} (Doc.~301, second level) and as such
necessarily incomplete; no description reconstructs the object in full. Fundamental in
the full sense is only the complete object (third level). That the resonance-faithful
geometry traces a \emph{real} object is therefore an explicit ontological \emph{posit}
--- not a derivation, and not itself falsifiable; what is falsifiable are only its
resonance-faithful consequences (the masses). This posit is named as such, not passed off as derived.

Fundamental is the information unit as a \emph{winding quantum} $\Delta w = 1$
(scale-universal, Doc.~257); \emph{energy}, by contrast, is a scale-specific
\emph{projection} of the geometry, $E_\text{bit}(L) = \hbar c / L$. An \emph{area}
shares this projected, scale-dependent character: the four-torus is boundaryless
($\partial T^4 = \emptyset$) and carries no horizon surface for bits to sit on; an area
there is a \emph{projection surface}, a second-level object. Where an entropy--area
relation $S = N\ln 2$ with $N \propto A/\xi^2$ appears, $N$ counts winding quanta and
$\ln 2$ is the thermodynamic shadow; the $A$ is not a fundamental partition of space
(which would revive Takayuki's cell worry) but a chosen projection. Fundamental (third
level) is neither area nor energy but the complete object; the scale-universal unit on
it is $\Delta w = 1$.

The honest accounting is thus: \emph{declared assumption} --- the compact geometry
($T^4/\mathbb{Z}_3$, $\xi$); \emph{definition} --- the identification \enquote{bit
$\widehat{=}\,\Delta w = 1$}, topologically grounded; \emph{derived} --- the
discreteness ($\pi_1(T^4)\cong\mathbb{Z}^4$), the scale-specific bit energy
$E_\text{bit} = \hbar c/L$, and the invariance; \emph{empirical exposure} --- the
masses; \emph{named posit} --- the reality of the object.

\section{Summary}

The objections are correct against a spatial-cell ontology, and the framework agrees
with them by not being one. The discreteness is topological ($\pi_1(T^4)\cong\mathbb{Z}^4$)
and spectral (the spectrum on a compact manifold); the minimal information unit is the
winding quantum $\Delta w = 1$ (the \enquote{bit}, by definition, topologically
grounded), not the capacity of a spatial cell; and the invariance holds because winding
numbers and spectra are invariants, not because a partition was chosen carefully.

\section{Where this is worked out}

The discrete spectrum of the $T^4$ connection-Laplacian is treated in Doc.~283; the
Hilbert-space formulation $\mathcal{H}=L^2(T^4)\otimes\mathbb{C}^3$ with masses as
eigenvalues in Doc.~282; the information unit as the winding quantum $\Delta w = 1$ and
the three-level hierarchy (topological / geometric / thermodynamic) in Doc.~257
(building on Doc.~255); the origin/orientation independence (\enquote{meshing, not a
cut-out}) in Doc.~300. The invariance point in particular deserves a fuller written
treatment than it has so far received.

